\section*{Chapitre \ref{ch:b1:rate-laws} --- Cinétique chimique~: lois de vitesse}

\begin{solution}{exo:b1:rate-laws:1}
$v = \frac14 \times 2.0 \times 10^{-4} = \qty{5.0e-5}{mol.L^{-1}.s^{-1}}$~;
le dioxygène se forme à la même vitesse, \qty{5.0e-5}{}~; \ce{N2O5}
disparaît à $2v = \qty{1.0e-4}{mol.L^{-1}.s^{-1}}$.
\end{solution}

\begin{solution}{exo:b1:rate-laws:2}
Ordre 0~: \unit{mol.L^{-1}.s^{-1}}~; ordre 1~: \unit{s^{-1}}~; ordre 2~:
\unit{L.mol^{-1}.s^{-1}}~; ordre 3~: \unit{L^2.mol^{-2}.s^{-1}}.
\end{solution}

\begin{solution}{exo:b1:rate-laws:3}
$t_{1/2} = \ln2/k = 0.693/(3.0 \times 10^{-3}) = \qty{231}{s}$. Au bout
de \qty{600}{s}, $\eu^{-kt} = \eu^{-1.8} = 0.165$~: il reste 16.5\,\%.
\end{solution}

\begin{solution}{exo:b1:rate-laws:4}
Un \omterm{def:b1:rate-laws:half-life}{temps de demi-réaction} proportionnel à $1/a_0$ signifie l'ordre 2. Si
diviser $a_0$ par deux divise par deux le \omterm{def:b1:rate-laws:half-life}{temps de demi-réaction},
$t_{1/2} \propto a_0$~: ordre 0.
\end{solution}

\begin{solution}{exo:b1:rate-laws:5}
$[\ce{A}]$ n'est pas affine en $t$ (pertes 0.0259, 0.0192, 0.0142,
0.0106). $\ln[\ce{A}] = -2.303$, $-2.602$, $-2.902$, $-3.202$, $-3.503$~:
il diminue de 0.300 toutes les \qty{10}{min}, c'est une droite. Ordre 1,
$k =
\qty{0.0300}{min^{-1}}$.
\end{solution}

\begin{solution}{exo:b1:rate-laws:6}
Doubler $[\ce{A}]_0$ multiplie $v_0$ par $4 = 2^2$~: ordre 2 par rapport
à \ce{A}. Doubler $[\ce{B}]_0$ la multiplie par 2~: ordre 1 par rapport
à \ce{B}. $v =
k[\ce{A}]^2[\ce{B}]$ avec $k = 2.0 \times 10^{-6}/(0.010^2 \times 0.010) =
\qty{2.0}{L^2.mol^{-2}.s^{-1}}$.
\end{solution}

\begin{solution}{exo:b1:rate-laws:7}
\ce{B} est en excès d'un facteur cent~: sa concentration reste égale à
\qty{0.50}{mol/L}, et $v = k'[\ce{A}]$ avec $k' = k[\ce{B}]_0$. D'où
$k = 0.012/0.50 = \qty{0.024}{L.mol^{-1}.s^{-1}}$.
\end{solution}

\begin{solution}{exo:b1:rate-laws:8}
$E_a = R\ln5/(1/300 - 1/320) = 8.314 \times 1.609/(2.083 \times 10^{-4}) =
\qty{64.2}{kJ/mol}$. À \qty{310}{K}~: $\ln k = \ln(1.0 \times 10^{-3}) +
\frac{E_a}{R}\big(\frac{1}{300} - \frac{1}{310}\big) = -6.908 + 0.831$,
$k = \qty{2.3e-3}{s^{-1}}$.
\end{solution}

\begin{solution}{exo:b1:rate-laws:9}
Ici $a = 2$~: $1/[\ce{A}] = 1/a_0 + 2kt$, $t_{1/2} = 1/(2k a_0) = 1/(2
\times 0.50 \times 0.020) = \qty{50}{s}$. À \qty{100}{s}~: $1/[\ce{A}] = 50
+ 100 = \qty{150}{L/mol}$, $[\ce{A}] = \qty{6.7e-3}{mol/L}$.
\end{solution}

\begin{solution}{exo:b1:rate-laws:10}
$[\ce{A}] = a_0(1 - f) = a_0\eu^{-akt_f}$ donne $t_f = -\ln(1-f)/(ak)$,
et, en \omterm{def:b1:rate-laws:half-life}{temps de demi-réaction}, $t_f/t_{1/2} = -\ln(1-f)/\ln2$~: 1 pour
50\,\%, 2 pour 75\,\%, 3.32 pour 90\,\%, 9.97 pour 99.9\,\%~: il faut
environ dix \omterm{def:b1:rate-laws:half-life}{temps de demi-réaction} pour que la réaction soit «~finie~».
\end{solution}

\begin{solution}{exo:b1:rate-laws:11}
$m(t) = 10 - 0.40\,t$ (\unit{mg}, \unit{h})~: ordre 0. \omterm{def:b1:rate-laws:half-life}{Temps de
demi-réaction} $10/(2 \times 0.40) = \qty{12.5}{h}$~; patch vide au bout
de \qty{25}{h}. La dose délivrée par heure ne dépend pas de ce qui
reste~: un apport régulier, ce qui est le but d'un patch.
\end{solution}

\begin{solution}{exo:b1:rate-laws:12}
$E_a = R\ln2/(1/298.15 - 1/308.15) = 8.314 \times 0.693/(1.088 \times
10^{-4}) = \qty{52.9}{kJ/mol}$. Entre 273.15 et \qty{283.15}{K}~:
$\exp\big(\frac{52900}{8.314}(\frac{1}{273.15} - \frac{1}{283.15})\big) =
2.3$~: la même \omterm{def:b1:rate-laws:activation-energy}{énergie d'activation} donne un facteur plus grand à plus
basse température.
\end{solution}

\begin{solution}{pb:b1:rate-laws:1}
\textbf{1.} Pertes de 8.6, 7.9, 7.2 et 6.5 points~: non constantes, donc
pas d'ordre 0.
\textbf{2.} $\ln$~: 4.605, 4.515, 4.425, 4.335, 4.245~; il diminue de
0.090 tous les 10 jours~: une droite, ordre 1.
\textbf{3.} $k_{60} = 0.090/10 = \qty{9.0e-3}{d^{-1}}$.
\textbf{4.} $t_{1/2} = \ln2/k_{60} = \qty{77}{d}$.
\textbf{5.} $100\,\eu^{-0.54} = 58.3\,\%$.
\textbf{6.} $t_{90} = \ln(10/9)/k_{60} = 0.1054/0.0090 = \qty{11.7}{d}$.
\textbf{7.} Il reste 90\,\%~: $\eu^{-kt_{90}} = 0.9$, donc
$t_{90} = \ln(10/9)/k$.
\textbf{8.} $t_{90}/t_{1/2} = \ln(10/9)/\ln2 = 0.152$.
\textbf{9.} Pour l'ordre 1, la fraction restante ne dépend que de $kt$,
et non de la quantité initiale.
\textbf{10.} À \qty{25}{\celsius}, la dégradation prend des années~; le
chauffage l'accélère (Arrhenius), de sorte qu'on peut la mesurer en
quelques semaines.
\textbf{11.} Pour l'ordre 2, le \omterm{def:b1:rate-laws:half-life}{temps de demi-réaction} vaut
$1/(ak\,a_0)$~: un comprimé deux fois plus dosé se dégraderait deux fois
plus vite en valeur relative et se conserverait moins longtemps.
\textbf{12.} $1/T$ (\unit{K^{-1}})~: \num{3.1934e-3}, \num{3.0945e-3},
\num{3.0017e-3}~; $\ln k$~: $-6.669$, $-5.661$, $-4.711$.
\textbf{13.} Pente $= (-4.711 + 6.669)/(3.0017 \times 10^{-3} - 3.1934 \times
10^{-3}) = \qty{-1.021e4}{K}$.
\textbf{14.} $E_a = 8.314 \times 1.021 \times 10^4 = \qty{85}{kJ/mol}$.
\textbf{15.} Valeur prévue $\ln k_{50} = -4.711 - 1.021 \times 10^4 \times (3.0945 -
3.0017) \times 10^{-3} = -5.658$, valeur mesurée $-5.661$~: le point est
sur la droite.
\textbf{16.} $A = k_{60}\,\eu^{E_a/RT} = 0.0090\,\eu^{10215/333.15} =
\qty{1.9e11}{d^{-1}}$.
\textbf{17.} $\exp\big(1.021 \times 10^4(\frac{1}{298.15} -
\frac{1}{308.15})\big) = 3.0$~: plus que le facteur 2 de la règle, parce
que $E_a$ dépasse \qty{53}{kJ/mol}.
\textbf{18.} $k_{30} = k_{60}\exp\big(-1.021 \times 10^4(\frac{1}{303.15} -
\frac{1}{333.15})\big) = \qty{4.3e-4}{d^{-1}}$.
\textbf{19.} $t_{90} = 0.1054/4.33 \times 10^{-4} = \qty{243}{d}$, soit
8.0 mois.
\textbf{20.} $1 - \eu^{-3 \times 0.0090} = 2.7\,\%$.
\textbf{21.} La durée de conservation diminue vite avec la température (8
mois à \qty{30}{\celsius} contre 14 à \qty{25}{\celsius})~: la date de
péremption n'est valable qu'en dessous de \qty{25}{\celsius}.
\textbf{22.} $k_{25} = k_{60}\exp\big(-1.021 \times 10^4(\frac{1}{298.15} -
\frac{1}{333.15})\big) = \qty{2.46e-4}{d^{-1}}$.
\textbf{23.} $t_{90} = 0.1054/2.46 \times 10^{-4} = \qty{428}{d}$, soit
\textbf{14 mois}~: la date imprimée sur la boîte, obtenue à partir de six
semaines de mesures.
\end{solution}
